\section*{Capítulo \ref{ch:b2:maxwell-equations} --- As equações de Maxwell}

\begin{solution}{exo:b2:maxwell-equations:1}
(a) $\operatorname{div} = 3$, $\operatorname{\vect{curl}} = \vect 0$: o \omterm{def:b2:maxwell-equations:operators}{gradiente} de $r^2/2$.
(b) $0$ e $(0, 0, 2)$: um \omterm{def:b2:maxwell-equations:operators}{rotacional} (de $-\tfrac12(x^2 + y^2)\vect e_z$), não um
\omterm{def:b2:maxwell-equations:operators}{gradiente}. (c) $0$ e $\vect 0$: o \omterm{def:b2:maxwell-equations:operators}{gradiente} de $xyz$ (e também um \omterm{def:b2:maxwell-equations:operators}{rotacional}). (d)
$\frac1{r^2}\partial_r(r^2/r^2) = 0$, \omterm{def:b2:maxwell-equations:operators}{rotacional} $\vect 0$: o \omterm{def:b2:maxwell-equations:operators}{gradiente} de $-1/r$. (e)
$0$ e $(0, -2x, 0)$: um \omterm{def:b2:maxwell-equations:operators}{rotacional}, não um \omterm{def:b2:maxwell-equations:operators}{gradiente}.
\end{solution}

\begin{solution}{exo:b2:maxwell-equations:2}
(a) Gauss e Faraday: os dois lados em \unit{V/m^2} (uma \omterm{def:b2:charges-currents-conduction:densities}{densidade de carga} sobre
$\varepsilon_0$, um \omterm{def:b2:maxwell-equations:operators}{gradiente} de campo, um campo sobre um tempo); Ampère: os dois lados em
\unit{T/m}.
(b) $1/\sqrt{8.85 \times 10^{-12} \times 1.257 \times 10^{-6}} = \qty{3.00e8}{m/s}$. (c) Razão
$\varepsilon_0\omega/\gamma$: cobre $5 \times 10^{-20}$; água do mar a \qty{1}{GHz} $1 \times 10^{-2}$
(ainda condutora, no limite); vidro a \qty{50}{Hz} $3 \times 10^3$: um dielétrico,
em que a \omterm{thm:b2:maxwell-equations:maxwell}{corrente de deslocamento} domina.
\end{solution}

\begin{solution}{exo:b2:maxwell-equations:3}
$\dd E/\dd t = I/\varepsilon_0\pi R^2 = \qty{1.8e12}{V.m^{-1}.s^{-1}}$; $j_d = I/\pi R^2 =
\qty{16}{A/m^2}$; $B(\qty{5}{cm}) = \mu_0Ir/2\pi R^2 = \qty{0.5}{\micro T}$; $B(\qty{20}{cm}) =
\mu_0I/2\pi r = \qty{0.5}{\micro T}$ --- exatamente o campo do fio.
\end{solution}

\begin{solution}{exo:b2:maxwell-equations:4}
(a) $\rho = \varepsilon_0k$. (b) $\frac1{r^2}\partial_r(\rho_0r^3/3\varepsilon_0) = \rho_0/\varepsilon_0$;
fora, por Gauss: $E = \rho_0R^3/3\varepsilon_0r^2$. (c) $\operatorname{div}\vect B = 0$: sim;
$\operatorname{\vect{curl}}\vect B = -(B_0/a)\vect e_z$: $\vect j = -(B_0/\mu_0a)\vect e_z$, uma
lâmina de corrente uniforme com espessura segundo $y$.
\end{solution}

\begin{solution}{exo:b2:maxwell-equations:5}
(a) $\vect A = \tfrac12(B_yz - B_zy, B_zx - B_xz, B_xy - B_yx)$: $(\operatorname{\vect{curl}}
\vect A)_x = \tfrac12(B_x + B_x) = B_x$, etc.; $\operatorname{div}\vect A = 0$. (b)
$\vect A' - \vect A = \tfrac12B_0(y, x, 0) = \operatorname{\vect{grad}}(\tfrac12B_0xy)$. (c)
$\frac1r\partial_r(Br^2/2) = B$; $\frac1r\partial_r(BR^2/2) = 0$; ambos valem $BR/2$ em
$R$. (d) $2\pi r\cdot BR^2/2r = \pi R^2B = \Phi$: uma espira fora do solenoide,
onde $\vect B = \vect 0$, ainda sente $e = -\dd\Phi/\dd t$ --- através de $\vect E =
-\partial_t\vect A$, que não se anula ali.
\end{solution}

\begin{solution}{exo:b2:maxwell-equations:6}
(a) $E = \sigma/\varepsilon_0 = \qty{1.1e5}{V/m}$, normal. (b) $\sigma/\varepsilon_0$ entre eles, nulo
fora (cada salto acrescenta $\pm\sigma/\varepsilon_0$). (c) $j_s = nI = \qty{1e4}{A/m}$, salto
$\mu_0j_s = \qty{12.6}{mT} = B$ no interior. (d) Uma espira fora do toro não enlaça
corrente resultante alguma, logo $B = 0$ ali; no interior, Ampère dá $\mu_0NI/2\pi r$,
e o salto através do enrolamento, $\mu_0j_s = \mu_0NI/2\pi r$, concorda.
\end{solution}

\begin{solution}{exo:b2:maxwell-equations:7}
(a) $V'' = -\rho_0/\varepsilon_0$: $V = -\rho_0x^2/2\varepsilon_0 + (U/d + \rho_0d/2\varepsilon_0)x$. (b)
$E = \rho_0x/\varepsilon_0 - U/d - \rho_0d/2\varepsilon_0$, nulo em $x_0 = d/2 + \varepsilon_0U/\rho_0d$
--- perto do meio para uma carga espacial grande. (c) $\sigma(0) = \varepsilon_0E(0^+) =
-\varepsilon_0U/d - \rho_0d/2$, $\sigma(d) = -\varepsilon_0E(d^-) = \varepsilon_0U/d - \rho_0d/2$: juntas,
$-\rho_0d$, neutralizando a carga espacial. (d) Com $\rho_0 = 0$: $V$ linear, $E$
uniforme; com $U = 0$: uma parábola com máximo em $d/2$.
\end{solution}

\begin{solution}{exo:b2:maxwell-equations:8}
(a) \qty{6000}{km}, \qty{300}{m}, \qty{3}{m}, \qty{12.5}{cm}. (b) $f < c/10L$:
\qty{300}{MHz}; \qty{3}{GHz}; \qty{30}{Hz}. (c) $(2\pi \times 10^6 \times 0.05/3 \times 10^8)^2/4 =
3 \times 10^{-7}$. (d) A rigor, não é: a \qty{50}{Hz} uma rede de \qty{1000}{km} vale um
sexto de comprimento de onda, as fases diferem ao longo dela e as linhas longas são
tratadas como \omterm{prop:b2:dispersion-wave-packets:coax}{linhas de transmissão} --- cada circuito local, porém, é
quase estacionário.
\end{solution}

\begin{solution}{exo:b2:maxwell-equations:9}
$\operatorname{div}\operatorname{\vect{curl}}\vect B = 0$ obrigaria $\operatorname{div}\vect j =
0$, ou seja, $\partial_t\rho = 0$ em toda parte. Na placa (volume $\pi R^2e$) a
corrente $I$ entra e nada sai: $\langle\operatorname{div}\vect j\rangle =
-I/\pi R^2e = \qty{-1.6e5}{A/m^3}$, $\partial_t\rho = \qty{+1.6e5}{C.m^{-3}.s^{-1}}$.
\end{solution}

\begin{solution}{exo:b2:maxwell-equations:10}
(a) $2\pi rE_\theta = -\pi r^2\dot B$: $E_\theta = -r\dot B/2$ no interior, $-R^2\dot B/2r$
fora. (b) Suas linhas são círculos fechados: ele não é um \omterm{def:b2:maxwell-equations:operators}{gradiente}, e
arrasta cargas ao longo de uma espira --- um campo eletromotor. (c) $e\cdot2\pi r_0
|E_\theta| = e\pi r_0^2\dot B$: $\pi \times 0.25 \times 10 = \qty{7.9}{eV}$ por volta; $2.5 \times 10^6$
voltas. (d) $p = eB(r_0)r_0$ dá $\dot p = er_0\dot B(r_0)$, ao passo que o campo
induzido dá $\dot p = e|E_\theta| = e\dot\Phi/2\pi r_0$: $B(r_0) = \Phi/2\pi r_0^2 =
\tfrac12\langle B\rangle$.
\end{solution}

\begin{solution}{exo:b2:maxwell-equations:11}
(a) $\operatorname{\vect{curl}}\operatorname{\vect{curl}}\vect E = -\Delta\vect E$ (pois $\operatorname{div}
\vect E = 0$) $= -\partial_t\operatorname{\vect{curl}}\vect B = -\mu_0\varepsilon_0\partial_t^2\vect E$. (b)
O mesmo com os papéis trocados. (c) $c = 1/\sqrt{\mu_0\varepsilon_0} = \qty{3.0e8}{m/s}$.
(d) Weber e Kohlrausch tinham medido a razão entre as unidades
eletrostática e eletromagnética de carga, $3.1 \times 10^8$ m/s; Fizeau havia
medido a luz a $3.15 \times 10^8$ m/s: "dificilmente podemos evitar a
conclusão de que a luz consiste em ondulações transversais do
mesmo meio".
\end{solution}

\begin{solution}{exo:b2:maxwell-equations:12}
(a) $\operatorname{\vect{curl}}\operatorname{\vect{curl}}\vect B = -\Delta\vect B = \mu_0\gamma\operatorname{\vect{curl}}
\vect E = -\mu_0\gamma\partial_t\vect B$: $\Delta\vect B = \mu_0\gamma\partial_t\vect B$, a forma da
equação do calor com difusividade $1/\mu_0\gamma$. (b) $\tau \sim \mu_0\gamma d^2$: \qty{7.5}{ms}
para a placa; $4\pi \times 10^{-7} \times 5 \times 10^5 \times 9 \times 10^{12} = \qty{6e12}{s}$, duzentos
mil anos, para o núcleo. (c) $\mu_0\gamma\delta^2 \sim 2/\omega$: $\delta = \sqrt{2/\mu_0
\gamma\omega}$: \qty{9}{mm} a \qty{50}{Hz}, \qty{65}{\micro m} a \qty{1}{MHz}. (d) O campo
do núcleo não pode mudar por difusão em menos de $10^5$ anos: a
variação secular observada e as inversões exigem o movimento do
fluido condutor --- um dínamo.
\end{solution}

\begin{solution}{pb:b2:maxwell-equations:1}
\textbf{1.} $Q = (I_0/\omega)\sin\omega t$, $E = Q/\varepsilon_0\pi R^2$: amplitude $I_0/\omega\varepsilon_0
\pi R^2 = \qty{5.7e5}{V/m}$.

\textbf{2.} $j_d = \varepsilon_0\partial_tE = I(t)/\pi R^2 = 32\cos\omega t$ A/m$^2$; seu fluxo
pela placa vale $I(t)$.

\textbf{3.} $2\pi rB = \mu_0\varepsilon_0\pi r^2\partial_tE$: $B = \mu_0I(t)r/2\pi R^2$; em $R$:
$\mu_0I_0/2\pi R = \qty{2}{\micro T}$.

\textbf{4.} Para $r > R$ toda a \omterm{thm:b2:maxwell-equations:maxwell}{corrente de deslocamento} fica enlaçada: $B =
\mu_0I/2\pi r$, o campo do fio; contínuo em $R$.

\textbf{5.} $(\operatorname{\vect{curl}}\vect E_1)_\theta = -\partial_rE_1 = -\partial_tB_\theta$, logo $\partial_rE_1
= \mu_0\varepsilon_0(r/2)\ddot E$ e $E_1(r) - E_1(0) = \mu_0\varepsilon_0r^2\ddot E/4$; com
$\ddot E = -\omega^2E$: $|E_1(R) - E_1(0)|/|E| = (\omega R/c)^2/4 = 1 \times 10^{-6}$.

\textbf{6.} Excelente. Para $10\%$ seria preciso $(\omega R/c)^2 = 0.4$, $f \approx \qty{300}{MHz}$ ---
uma cavidade, e não mais um capacitor.

\textbf{7.} $W_e = \tfrac12\varepsilon_0E_0^2\pi R^2d = \qty{4.5e-5}{J}$; $W_m = \int_0^R(\mu_0I_0r/
2\pi R^2)^2/2\mu_0\cdot2\pi rd\,\dd r = \mu_0I_0^2d/16\pi = \qty{2.5e-11}{J}$; razão $6 \times
10^{-7} \approx (\omega R/c)^2/8$: o mesmo pequeno parâmetro.

\textbf{8.} $\oint\vect E\cdot\dd\vect l = 2\pi r_0E_\theta = -\dd\Phi/\dd t$.

\textbf{9.} $\dd p/\dd t = e|E_\theta| = (e/2\pi r_0)|\dd\Phi/\dd t|$; por volta,
$2\pi r_0\cdot e|E_\theta| = e|\dd\Phi/\dd t|$.

\textbf{10.} $p = eB(r_0)r_0$ dá $\dot p = er_0\dot B(r_0)$; igualando com 9:
$B(r_0) = \Phi/2\pi r_0^2 = \tfrac12\langle B\rangle$.

\textbf{11.} Um campo uniforme tem $B(r_0) = \langle B\rangle$, o dobro do necessário:
a força magnética cresce mais depressa que a quantidade de movimento e a órbita
encolhe. As peças polares são moldadas de modo que o campo seja forte perto do
eixo e metade disso na órbita.

\textbf{12.} $p = eBr_0 = \qty{3.2e-20}{kg.m/s}$; $E = pc = \qty{9.6e-12}{J} = \qty{60}{MeV}$;
$\Phi = 2\pi r_0^2B = \qty{0.63}{Wb}$ em \qty{5}{ms}: \qty{126}{eV} por volta, $4.8 \times
10^5$ voltas, \qty{1500}{km} --- a $c$ em \qty{5}{ms}: coerente.

\textbf{13.} Da fonte de alimentação do ímã, entregue aos elétrons pelo
termo de Faraday: o fluxo variável cria o campo acelerador.

\textbf{14.} Dezenas de MeV de elétrons freando num alvo pesado dão
raios X penetrantes; o acelerador linear compacto, de taxas de dose
mais altas, o substituiu.

\textbf{15.} $\operatorname{div}\operatorname{\vect{curl}}\vect A = 0$; $\operatorname{\vect{curl}}(-\operatorname{\vect{grad}}
V - \partial_t\vect A) = -\partial_t\operatorname{\vect{curl}}\vect A = -\partial_t\vect B$.

\textbf{16.} $\frac1{r^2}(r^2V')' = -\rho_0/\varepsilon_0$: $V = \rho_0(3a^2 - r^2)/6\varepsilon_0$
no interior, $\rho_0a^3/3\varepsilon_0r = Q/4\pi\varepsilon_0r$ fora; $E = \rho_0r/3\varepsilon_0$ e
$Q/4\pi\varepsilon_0r^2$.

\textbf{17.} $\tfrac12\int\rho V\,\dd\tau = 2\pi\rho_0^2(a^5/6 - a^5/30)/\varepsilon_0 = 4\pi\rho_0^2a^5/
15\varepsilon_0 = 3Q^2/20\pi\varepsilon_0a$; urânio: \qty{1.6e-10}{J} $\approx \qty{1000}{MeV}$ (a fissão
libera os \qty{200}{MeV} de diferença por que duas bolas menores saem mais baratas).

\textbf{18.} $V = -E(t)z$ (a menos de uma constante). Com $\vect A = A_z(r, t)\vect e_z$,
$B_\theta = -\partial_rA_z = \mu_0Ir/2\pi R^2$ dá $A_z = -\mu_0I(t)r^2/4\pi R^2$.

\textbf{19.} $-\partial_tA_z = \mu_0\dot Ir^2/4\pi R^2 = \mu_0\varepsilon_0r^2\ddot E/4$ (pois $\dot I =
\varepsilon_0\pi R^2\ddot E$): a correção da questão 5, a menos de uma constante.

\textbf{20.} $\operatorname{div}\vect A = \partial_zA_z = 0$ porque $A_z$ só depende de $r$;
qualquer $\vect A + \operatorname{\vect{grad}}\chi$ serve.

\textbf{21.} $\sigma = \varepsilon_0E_0 = \qty{5}{\micro C/m^2}$; dentro do metal $\vect E =
\vect 0$, de modo que o salto $\sigma/\varepsilon_0$ é exatamente o campo do vão.

\textbf{22.} A corrente $I(t)$ alimenta a placa a partir do eixo; a
parte além de $r$ que ainda falta entregar vale $I(1 - r^2/R^2)$, logo $j_s = I(1 -
r^2/R^2)/2\pi r$. Na face externa Ampère dá o campo do fio $\mu_0I/
2\pi r$; na face interna, $\mu_0Ir/2\pi R^2$; a diferença, $\mu_0I(1 -
r^2/R^2)/2\pi r = \mu_0j_s$, é o salto de passagem.

\textbf{23.} Placas: $R/\lambda = 3 \times 10^{-4}$, quase estacionário elétrico. Espira: $7 \times 10^{-4}$,
quase estacionário magnético. A \qty{1}{GHz}: $R/\lambda = 0.3$ --- uma cavidade e uma antena.

\textbf{24.} $\lambda = \qty{3000}{km}$ para um anel de alguns metros: a \omterm{thm:b2:maxwell-equations:maxwell}{corrente de
deslocamento} é totalmente desprezível; o campo elétrico essencial é o
de Faraday, que o \omterm{def:b2:maxwell-equations:arqs}{regime quase estacionário} magnético mantém por inteiro.

\textbf{25.} I: Ampère--Maxwell (\omterm{thm:b2:maxwell-equations:maxwell}{corrente de deslocamento}), Faraday para
a correção, densidades de energia. II: Faraday integral, força de Lorentz.
III: os potenciais, Poisson (Gauss). IV: os saltos de passagem de $\vect E$
e $\vect B$; a razão tamanho/comprimento de onda para o regime.
\end{solution}
